% Second numerical example: flow past a centered square cylinder.
% Recommended packages in the main manuscript:
% \usepackage{amsmath,amssymb,booktabs,graphicx,threeparttable}
% In a single-column manuscript, table* and figure* may be replaced by
% table and figure, respectively.

\subsection{Flow past a centered square cylinder}
\label{subsec:flow_past_centered_square_cylinder}

The flow past a centered square cylinder is used to assess RAL-MoE-ROM for a
different geometry and spatial discretization. The numerical study examines
the long-horizon accuracy of the regime-local specialists and the T2-C fusion
in both the $S$--$H$ and $H$--$P$ overlap regions.

The high-fidelity model solves the two-dimensional incompressible
Navier--Stokes equations in the channel domain
$\Omega=[0,31D]\times[0,4D]$. A square cylinder of side length $D=1$
occupies $[5D,6D]\times[1.5D,2.5D]$, corresponding to a blockage ratio of
$0.25$. The inlet velocity profile is
\begin{equation}
u_x(y)
=
4U_m\frac{y(4D-y)}{(4D)^2},
\qquad
u_y=0,
\qquad
U_m=1,
\label{eq:centered_square_inlet}
\end{equation}
with zero-normal-gradient velocity at the outlet and no-slip conditions on
the channel walls and square cylinder. Pressure is fixed at the outlet and
satisfies zero-normal-gradient conditions on the remaining boundaries. The
Reynolds number is
\[
Re=\frac{U_mD}{\nu}.
\]


The full-order simulations are performed with OpenFOAM~13 using
\texttt{icoFoam} and Crank--Nicolson time integration. The graded
finite-volume mesh contains 9,400 cells, with refinement around the obstacle
and in the near wake. The CFD time step is
$\delta t_{\mathrm{CFD}}=0.02$, and snapshots are stored every 200 CFD
steps, giving
\begin{equation}
\Delta t_{\mathrm{snap}}=4.
\label{eq:centered_square_snapshot_dt}
\end{equation}
The parameterized dataset contains 120 Reynolds-number trajectories over
$50\le Re\le150$, with additional sampling in the transition region.


Table~\ref{tab:centered_square_data_contracts} summarizes the three
overlapping regime-local datasets, their complete-trajectory partitions,
local POD dimensions, and rollout horizons. All regime-local means, POD
bases, and normalization statistics are constructed exclusively from the
corresponding training trajectories. The velocity and pressure reduced
dimensions are selected to retain $99.9\%$ of the cumulative POD energy.
All three specialists follow the common formulation introduced in
Section~3.

\begin{table*}[htbp]
\centering
\caption{Data partitions, local POD dimensions, and rollout horizons for flow past a centered square cylinder.}
\label{tab:centered_square_data_contracts}

\scriptsize
\begin{adjustbox}{max width=\textwidth}
\begin{tabular}{@{}lccccccc@{}}
\toprule
Regime
& Reynolds-number range
& Train/validation/test
& $r_u$
& $r_p$
& Rollout horizon
& Physical-time horizon \\
\midrule

Steady
& $50\le Re\le95.40$
& $60/5/4$
& 5
& 4
& $K=16$
& 64 \\

Hopf
& $94\le Re\le102$
& $23/6/5$
& 11
& 11
& $K=48$
& 192 \\

Periodic
& $98.5\le Re\le150$
& $27/6/4$
& 28
& 26
& $K=24$
& 96 \\

\bottomrule
\end{tabular}
\end{adjustbox}
\end{table*}

All model and fusion parameters used in the following held-out evaluations
are selected using training and validation data only and are fixed before
access to the corresponding test trajectories. Since
$\Delta t_{\mathrm{snap}}=4$ for all three datasets, the $K=16$, $K=48$,
and $K=24$ rollouts span 64, 192, and 96 physical time units,
respectively.

Table~\ref{tab:centered_square_specialist_aggregate} reports the
held-out long-horizon physical-field errors of the three proposed
regime-local specialists. These values are evaluated under their respective
POD representations and rollout horizons and are therefore used to assess
the performance of each specialist within its own regime rather than to form
a direct ranking across regimes.

\begin{table}[htbp]
\centering
\caption{Held-out long-horizon prediction errors of the regime-local
specialists for the centered-square-cylinder case.}
\vspace{0.3 cm}
\label{tab:centered_square_specialist_aggregate}

\small
\renewcommand{\arraystretch}{1.15}
\setlength{\tabcolsep}{12pt}

\begin{tabular}{@{}lccc@{}}
\toprule
Regime
& Rollout horizon
& $E_u$ (\%)
& $E_p$ (\%) \\
\midrule
Steady
& $K=16$
& 1.5311
& 1.4583 \\

Hopf
& $K=48$
& 0.4387
& 0.4844 \\

Periodic
& $K=24$
& 0.9395
& 1.0921 \\
\bottomrule
\end{tabular}
\end{table}

The three specialists maintain low physical-field errors over their
respective held-out long-horizon evaluations. The Steady specialist yields
velocity and pressure errors of $1.5311\%$ and $1.4583\%$ over $K=16$,
whereas the Hopf specialist achieves $0.4387\%$ and $0.4844\%$ over the
substantially longer $K=48$ horizon. For the Periodic specialist, the
corresponding $K=24$ errors are $0.9395\%$ and $1.0921\%$. These results
show that the regime-local construction remains accurate under the
square-cylinder geometry despite the markedly different POD dimensions and
dynamical regimes.


Having established the predictive accuracy of the three regime-local
specialists, we next examine the hierarchical coupling in the two parameter
regions where adjacent local models overlap. Both the $S$--$H$ and $H$--$P$
comparisons are performed over a common $K=24$ prediction horizon. To assess the contribution of the T2-C correction, the two overlap regions
are evaluated against three non-learned routing or fusion strategies and an
a posteriori convex reference, in addition to the individual regime-local
specialists. For an admissible adjacent pair, E2 Top-1 selects the specialist
with the larger pairwise E2 probability and uses its reconstructed field
without fusion. The equal-weight blend combines the two reconstructed fields
with fixed weights $1/2$ and $1/2$, whereas the E2-probability blend uses
the renormalized pair probabilities $\bar{\pi}_1$ and $\bar{\pi}_2$ from
Eq.~\eqref{eq:t2c_pair_logit} directly as the convex weights. T2-C employs
the learned window-conditioned correction defined in
Eq.~\eqref{eq:t2c_convex_gate}. A window-wise convex oracle is additionally
reported as an a posteriori reference by selecting, for each evaluation
window, the convex weight that minimizes the prescribed field-error
functional over $\alpha\in[0,1]$. Since the oracle uses the reference
solution, it is used only to quantify the attainable accuracy of convex
output combination and is not available during prediction.

\begin{table*}[htbp]
\centering
\caption{Prediction errors for the routing and fusion.xi}
\label{tab:centered_square_fusion_aggregate}

\small
\renewcommand{\arraystretch}{1.10}
\setlength{\tabcolsep}{8pt}

\begin{tabular}{@{}llccc@{}}
\toprule
Overlap
& Model
& $E_u$ (\%)
& $E_p$ (\%)
& $E_{\mathrm{joint}}$ (\%) \\
\midrule

$S$--$H$
& Steady specialist
& 2.5976
& 2.2450
& 4.8426 \\

& Hopf specialist
& 0.4942
& 1.2356
& 1.7298 \\

& E2 Top-1
& 0.4942
& 1.2356
& 1.7298 \\

& Equal-weight blend
& 1.3644
& 1.6677
& 3.0321 \\

& E2-probability blend
& 0.7898
& 1.3949
& 2.1847 \\

& T2-C
& 0.4615
& 0.7343
& 1.1958 \\

& Window-wise convex oracle
& 0.4659
& 0.7183
& 1.1842 \\

\midrule

$H$--$P$
& Periodic specialist
& 0.6420
& 0.5659
& 1.2079 \\

& Hopf specialist
& 1.0034
& 1.8684
& 2.8718 \\

& E2 Top-1
& 0.7422
& 1.1052
& 1.8473 \\

& Equal-weight blend
& 0.6077
& 1.0069
& 1.6146 \\

& E2-probability blend
& 0.6154
& 1.0464
& 1.6618 \\

& T2-C
& 0.4509
& 0.3956
& 0.8465 \\

& Window-wise convex oracle
& 0.4479
& 0.3905
& 0.8384 \\

\bottomrule
\end{tabular}
\end{table*}

The two overlap regions exhibit distinct patterns of predictive
complementarity between adjacent regime-local specialists. In the
$S$--$H$ overlap, E2 selects the Hopf specialist for both held-out
Reynolds-number trajectories, so that the aggregate E2 Top-1 errors coincide
with those of the Hopf specialist. The corresponding joint error is
$1.7298\%$. Neither the equal-weight blend nor the direct use of the
renormalized E2 probabilities as convex weights improves upon this result;
their joint errors increase to $3.0321\%$ and $2.1847\%$, respectively.
T2-C instead reduces $E_u$ from $0.4942\%$ to $0.4615\%$ and $E_p$ from
$1.2356\%$ to $0.7343\%$, yielding
$E_{\mathrm{joint}}=1.1958\%$. This corresponds to a $30.87\%$ reduction
relative to E2 Top-1 and is close to the window-wise convex oracle value of
$1.1842\%$. The improvement therefore arises from the window-conditioned
correction of the pairwise E2 weighting rather than from convex averaging
alone.

The routing behaviour is different in the $H$--$P$ overlap, where E2 selects
different specialists across the held-out Reynolds-number trajectories.
Consequently, the aggregate E2 Top-1 errors do not coincide with those of
either fixed individual specialist. Among the two single-specialist
predictions, the Periodic specialist gives the lower aggregate joint error,
$E_{\mathrm{joint}}=1.2079\%$, whereas E2 Top-1 yields
$E_{\mathrm{joint}}=1.8473\%$. The equal-weight and E2-probability blends
improve upon the hard-routing result, but remain less accurate than the
Periodic specialist alone. T2-C further reduces the velocity and pressure
errors to $0.4509\%$ and $0.3956\%$, respectively, giving
$E_{\mathrm{joint}}=0.8465\%$. This corresponds to reductions of $29.92\%$
relative to the best single-specialist baseline and $54.18\%$ relative to
E2 Top-1. The result is again close to the window-wise convex oracle
($0.8384\%$), indicating that the learned correction recovers most of the
attainable improvement within the considered class of output-level convex
combinations.

To examine whether these aggregate improvements are also observed across
the individual held-out Reynolds-number trajectories,
Table~\ref{tab:centered_square_fusion_parameter} reports the trajectory-level
E2 selections, the mean T2-C weights, and the corresponding joint errors.
The reported weight is $\bar{\alpha}_S$ for the $S$--$H$ overlap and
$\bar{\alpha}_P$ for the $H$--$P$ overlap.

\begin{table*}[htbp]
\centering
\caption{Parameter-wise routing and T2-C fusion results at $K=24$.
$\bar{\alpha}$ denotes the mean Steady weight in the $S$--$H$ overlap and
the mean Periodic weight in the $H$--$P$ overlap.}
\label{tab:centered_square_fusion_parameter}

\small
\renewcommand{\arraystretch}{1.10}
\setlength{\tabcolsep}{7pt}

\begin{tabular}{@{}lr l c ccc@{}}
\toprule
Overlap
& $Re$
& E2-selected specialist
& $\bar{\alpha}$
& $E_{\mathrm{joint}}^{\mathrm{E2}}$ (\%)
& $E_{\mathrm{joint}}^{\mathrm{T2-C}}$ (\%)
& Reduction (\%) \\
\midrule

$S$--$H$
& 95.1
& Hopf
& 0.0967
& 1.7140
& 1.1827
& 31.00 \\

$S$--$H$
& 95.3
& Hopf
& 0.0971
& 1.7456
& 1.2090
& 30.74 \\

$H$--$P$
& 100.5
& Hopf
& 0.5870
& 2.6485
& 0.8571
& 67.64 \\

$H$--$P$
& 102
& Periodic
& 0.7094
& 1.0462
& 0.8360
& 20.09 \\

\bottomrule
\end{tabular}
\end{table*}

For the $S$--$H$ overlap, the trajectory-level routing and fusion behaviour
is consistent across the two held-out Reynolds-number trajectories. E2
selects the Hopf specialist in both cases, while T2-C assigns only a small
mean weight to the Steady specialist,
$\bar{\alpha}_S\approx0.097$. Despite this predominantly Hopf-based
combination, T2-C reduces the joint error by approximately $31\%$ for both
trajectories. The improvement is therefore consistently observed across the
held-out $S$--$H$ overlap rather than being dominated by a single parameter
value.

The $H$--$P$ overlap exhibits a more pronounced variation in both the
trajectory-level E2 decisions and the T2-C weights. E2 selects different
specialists across the held-out trajectories, whereas the mean Periodic
weight produced by T2-C varies from $0.5870$ to $0.7094$. The corresponding
joint-error reductions range from $20.09\%$ to $67.64\%$, with improvements
obtained for both trajectories. These results indicate that the T2-C
correction does not reduce to a fixed preference for either adjacent
specialist; instead, the relative contribution of the two regime-local
predictions varies across the overlap.

Taken together, the two overlap studies separate the role of regime
selection from that of local-model combination. E2 provides a discrete
trajectory-level choice among admissible specialists, whereas T2-C exploits
the remaining predictive complementarity between adjacent local ROMs after
their independent evolution. The equal-weight and E2-probability baselines
show that this improvement is not reproduced by an arbitrary convex
combination, while the close agreement with the window-wise convex oracle
indicates that the learned correction approaches the best performance
available within the considered output-level convex fusion class.
The spatial distributions of the overlap predictions are shown in
Figs.~\ref{fig:centered_square_sh_fields_errors} and
\ref{fig:centered_square_hp_fields_errors}. These visualizations complement
the integral error measures by identifying the regions in which the adjacent
regime-local predictions differ and where the T2-C combination modifies the
reconstructed physical fields. 
\begin{figure*}[htbp]
\centering

\includegraphics[width=0.90\textwidth]
{DDPDCNN/4/RAL_MOE_ROM_METHODS_OVERLEAF_READY_20260726/drafts/cylinder_wake_20260726/figures/SH_2.png}

\vspace{0.5em}

\includegraphics[width=0.90\textwidth]
{DDPDCNN/4/RAL_MOE_ROM_METHODS_OVERLEAF_READY_20260726/drafts/cylinder_wake_20260726/figures/SH_2_1.png}

\caption{Representative held-out prediction in the $S$--$H$ overlap at
$Re=95.1$. The upper panel compares the CFD reference, the Steady
specialist, the Hopf specialist, and the T2-C prediction in terms of
velocity magnitude and gauge-corrected pressure. The lower panel shows the
corresponding pointwise absolute errors in velocity and pressure.}
\label{fig:centered_square_sh_fields_errors}
\end{figure*}

\begin{figure*}[htbp]
\centering

\includegraphics[width=0.90\textwidth]
{DDPDCNN/4/RAL_MOE_ROM_METHODS_OVERLEAF_READY_20260726/drafts/cylinder_wake_20260726/figures/HP_2.png}

\vspace{0.5em}

\includegraphics[width=0.90\textwidth]
{DDPDCNN/4/RAL_MOE_ROM_METHODS_OVERLEAF_READY_20260726/drafts/cylinder_wake_20260726/figures/HP_2_1.png}

\caption{Representative held-out prediction in the $H$--$P$ overlap at
$Re=100.5$. The upper panel compares the CFD reference, the Periodic
specialist, the Hopf specialist, and the T2-C prediction in terms of
velocity magnitude and gauge-corrected pressure. The lower panel shows the
corresponding pointwise absolute errors in velocity and pressure.}
\label{fig:centered_square_hp_fields_errors}
\end{figure*}

